%% ===================================================================
%% Paper 20 drop-in proposal — Tier 2 resource table
%% (Track T6, Dirac-on-S³ Tier 2 sprint).
%%
%% Target file:  papers/applications/paper_20_resource_benchmarks.tex
%% Target site:  appended to the resource-table section, as
%%               "Table 2" (after the existing scalar-composed
%%               benchmark table).
%% PI decision:  T6 proposes; PI applies at sprint commit.
%% ===================================================================

\subsection{Relativistic (spin-ful) composed resource benchmarks}
\label{subsec:paper20_tier2}

Table~\ref{tab:paper20_tier2} reports resource metrics for the
spin-ful composed encoding of Paper~14~\S\ref{sec:spinor_composed}
(Dirac-on-$S^3$ Tier 2 construction). Three single-bond hydrides
spanning $Z \in \{3, 4, 20\}$, at $n_{\max} = 2$ and $n_{\max} = 3$.
Every metric is computed from the algebraic-first pipeline described
in that section: exact-rational $X_k$ angular coefficients (sympy
\texttt{wigner\_3j}), exact-rational or machine-precision $R^k$ radial
integrals (\texttt{hypergeometric\_slater.py}), closed-form
$H_{\mathrm{SO}} = -Z^4\alpha^2(\kappa+1)/[4 n^3 l(l+\tfrac{1}{2})(l+1)]$
Breit--Pauli diagonal. No numerical quadrature enters the spinor path.

\begin{table*}[h]
\caption{Tier 2 relativistic composed resource metrics (Dirac-on-$S^3$
T3 output). $N_{\mathrm{Pauli}}^{\mathrm{rel}}$: non-identity Pauli
string count in the $(\kappa, m_j)$-labeled JW encoding. $\lambda_{\mathrm{ni}}$:
non-identity Hamiltonian 1-norm. QWC: greedy qubit-wise-commuting
group count. ``scalar'' repeats the spinless benchmark at the same
molecular spec topology. Sunaga ratio = $N_{\mathrm{Pauli}}^{\mathrm{rel}} / 47\,099$,
the single calibrated Sunaga RaH-18q cell~\cite{Sunaga2025}.}
\label{tab:paper20_tier2}
\begin{ruledtabular}
\begin{tabular}{cccccccc}
Molecule & $n_{\max}$ & $Q$ & $N_{\mathrm{Pauli}}^{\mathrm{rel}}$ & $N_{\mathrm{Pauli}}^{\mathrm{sc}}$ & $\lambda_{\mathrm{ni}}$ (Ha) & QWC & Sunaga ratio \\
\hline
LiH & 2 & 30 &    805  & 333    &  35.90 &   55 & 0.017$\times$ \\
LiH & 3 & 84 & 46\,434 & 7\,878 & 126.75 & 6\,571 & 0.986$\times$ \\
BeH & 2 & 30 &    805  & 333    & 141.32 &   52 & 0.017$\times$ \\
BeH & 3 & 84 & 46\,434 & 7\,878 & 297.03 & 6\,571 & 0.986$\times$ \\
CaH & 2 & 20 &    534  & 222    &  13.87 &   52 & 0.011$\times$ \\
CaH & 3 & 56 & 30\,940 & 5\,252 &  65.81 & 6\,571 & 0.657$\times$ \\
\hline\hline
\multicolumn{8}{l}{\footnotesize Sunaga RaH-18q~\cite{Sunaga2025}: 47\,099 Pauli (rel), 12\,556 Pauli (non-rel), $Q = 18$.} \\
\multicolumn{8}{l}{\footnotesize BeH / MgH / CaH / SrH / BaH at 18q: published in Sunaga SI Tables S1--S3, flagged \textbf{DEFERRED}.} \\
\end{tabular}
\end{ruledtabular}
\end{table*}

Reading Table~\ref{tab:paper20_tier2}:

\begin{itemize}
  \item At $n_{\max}=2$, all three GeoVac molecules sit at roughly
    $0.01$--$0.02\times$ the Sunaga RaH-18q Pauli count at native
    $Q \in \{20, 30\}$. Projected to matched $Q = 18$ via the
    O($Q^{2.5}$) scaling law of Paper~14~\S\ref{sec:composed} and the
    Tier 2 $\sim 2.4\times$ rel/scalar multiplier, the advantage is
    $150\times$--$250\times$.
  \item At $n_{\max}=3$, the absolute GeoVac Pauli count catches up to
    Sunaga at matched $Q$ (LiH/BeH approach 1$\times$ of the RaH-18q
    number at $Q = 84$; CaH sits at $0.66\times$ at $Q = 56$). This
    is the expected sparsity-exponent regime: as $l_{\max}$ grows,
    the composed architecture's per-block sparsity advantage narrows.
  \item The $\lambda_{\mathrm{ni}}$ column is non-inflated by
    relativistic coupling. For LiH $n_{\max}=3$, scalar gives
    170.59 Ha and rel gives 126.75 Ha ($-26\%$); for CaH,
    96.45 Ha $\to$ 65.81 Ha ($-32\%$). This is structurally favorable
    for QPE resource estimation (query complexity $\propto \lambda$).
  \item The QWC column is inflated ($\sim 8\times$ in the worst case).
    The richer $jj$-coupled Pauli structure distributes the
    Hamiltonian across more distinct Pauli words; VQE measurement
    overhead grows accordingly. Mitigable by symmetry-compressed
    or tensor-factorized measurement strategies.
\end{itemize}

\textbf{Deferred Sunaga cells.} The per-molecule Pauli counts, 1-norm,
and QWC values for BeH, MgH, CaH, SrH, BaH at $Q = 18$ are reported
in Sunaga et al.~\cite{Sunaga2025} Supplemental Material Tables
S1--S3. These were not extracted in the Tier 2 sprint timeframe. The
footnote row in Table~\ref{tab:paper20_tier2} marks this as deferred;
sharpening the head-to-head at matched $Q = 18$ per-molecule is
future work, noted in Paper~14~\S\ref{subsec:spinor_scope}.

\textbf{Accuracy caveats.} The Tier 2 spin-ful encoding inherits the
composed architecture's $5$--$26\,\%$ $R_{\mathrm{eq}}$ error from
Paper~17 and the Breit--Pauli leading-order $Z\alpha^2$ approximation
from Paper~14~\S\ref{sec:spinor_composed}. Sunaga's DC-UCCSD is
accuracy-targeted ($0.47\%$ vs CASCI for CaH). Table~\ref{tab:paper20_tier2}
is a \emph{resource estimation} table at fixed single-point
geometries, not a spectroscopic accuracy benchmark.

\textbf{Fine-structure sanity check.} The closed-form Breit-Pauli
$H_{\mathrm{SO}}$ reproduces the He ($2^3$P), Li ($2^2$P), and Be
($2s2p\,^3$P) fine-structure splittings to correct sign and correct
order of magnitude (benchmark code
\texttt{benchmarks/fine\_structure\_check.py}). Relative errors
$-66\%, +211\%, -78\%$ against NIST~\cite{NIST_ASD} --- consistent
with the leading-order $Z\alpha^2$ approximation and simple
$Z_{\mathrm{eff}}$ screening, not spectroscopic. The 20--50\%
accuracy target suggested in the Tier 2 sprint plan is not met; full
Dirac-Coulomb radial corrections and multi-electron Breit-Pauli
corrections are Tier 3+.
